\clearpage
\chapter*{Errata da exposição n.º 212}
\addcontentsline{toc}{chapter}{Errata da exposição n.º 212}
\setcounter{footnote}{2}
\markboth{Errata da exposição 212}{Errata da exposição 212}
\sid{e212}
\src{301}

\noindent\textbf{5. TDTE III: Pré-esquemas quocientes.}

\begin{center}
[Seminário Bourbaki, t. 13, 1960/61, n.º 212, 20 p.]
\end{center}

\noindent\textbf{Página 15, asserção (i).} --- O fato de que $Y=X/R$ seja
quase-projetivo sobre $S$ parece, por ora, ter sido demonstrado apenas quando
$R$ provém de uma relação de equivalência.

\noindent\textbf{Páginas 18 e 19.} --- Como assinalamos ao final da exposição
seguinte, a conjectura considerada é decididamente falsa. O «fato positivo»
assinalado na Observação 8.1 parece ter sido demonstrado simultaneamente por
diversos autores (NAGATA, ROSENLICHT, GROTHENDIECK, \ldots).

Assinalemos que a aplicação feita em TDTE V da teoria da passagem ao quociente
aqui desenvolvida à construção dos esquemas de Picard também pode ser
substituída por uma utilização adequada dos esquemas de Hilbert (cf. SHMT).
Como foi dito no § 8, a lacuna mais importante da teoria apresentada
\src{302}
aqui é a ausência de um critério de existência de quocientes por uma relação de
equivalência não própria, tais como as relações de equivalência provenientes de
certas ações do grupo projetivo. Um teorema importante nessa direção foi obtido
por MUMFORD\footnote{MUMFORD (David). --- An elementary theorem in geometric
invariant theory, \emph{Bull. Amer. Math. Soc.}, t. 67, 1961, pp. 483--486.}.
Para um refinamento de seu resultado e diversas aplicações à teoria, ver SHMT.

\par
